\documentclass[12pt, a4paper]{ctexart}
\usepackage{fancyhdr}
\usepackage{amsmath, amssymb}
\usepackage{geometry}
\usepackage{enumitem}
\usepackage{setspace}
\usepackage{graphicx}
\usepackage{tabularx}
\usepackage{array}
\usepackage{multirow}

\geometry{left=2cm, right=2cm, top=2.5cm, bottom=2.5cm}
\setCJKmainfont{Hiragino Sans GB}
\setCJKsansfont{Heiti SC}
\setstretch{1.5}

\pagestyle{fancy}
\fancyhf{}
\fancyhead[L]{}
\fancyhead[R]{\small 数学（选择性必修第三册）}
\fancyfoot[C]{\thepage}
\renewcommand{\headrulewidth}{0.4pt}

\newcommand{\score}[1]{\hfill\textbf{（#1分）}}

\begin{document}

\begin{center}
\Large\textbf{选择性必修第三册·第八章 成对数据的统计分析 测试卷}

\vspace{0.3cm}
\normalsize 满分：150分 \quad 考试时间：120分钟
\end{center}

\vspace{0.5cm}

\noindent\textbf{一、单项选择题（本大题共8小题，每小题5分，共40分。每小题只有一个选项符合题目要求。）}

\vspace{0.3cm}

\begin{enumerate}[label=\textbf{\arabic*.}, leftmargin=2em]

\item 变量$X$与$Y$的相关系数$r$的取值范围是\score{5}
\begin{enumerate}[label=\Alph*., leftmargin=2em]
\item $0 \leqslant |r| \leqslant 1$ \quad B. $-1 \leqslant r \leqslant 0$ \quad C. $0 \leqslant r \leqslant 1$ \quad D. $-1 \leqslant |r| \leqslant 1$
\end{enumerate}

\item 已知回归直线方程为$\hat{y} = 2x - 1$，当$x = 5$时，$y$的预测值为\score{5}
\begin{enumerate}[label=\Alph*., leftmargin=2em]
\item $7$ \quad B. $9$ \quad C. $11$ \quad D. $13$
\end{enumerate}

\item 在回归分析中，残差是指\score{5}
\begin{enumerate}[label=\Alph*., leftmargin=2em]
\item 观测值$y_i$与预测值$\hat{y}_i$的差
\item 观测值$y_i$与均值$\bar{y}$的差
\item 预测值$\hat{y}_i$与均值$\bar{y}$的差
\item 以上都不对
\end{enumerate}

\item 对两个变量$X$和$Y$进行线性回归分析，若相关系数$r = 0.95$，则下列说法正确的是\score{5}
\begin{enumerate}[label=\Alph*., leftmargin=2em]
\item $X$与$Y$之间具有确定的函数关系
\item $X$与$Y$之间具有很强的线性相关关系
\item $X$与$Y$之间不具有相关关系
\item $X$与$Y$之间具有因果关系
\end{enumerate}

\item 某同学根据一组数据$(x_i, y_i)$（$i=1,2,\ldots,10$）求得回归直线方程为$\hat{y} = \hat{b}x + \hat{a}$，已知$\bar{x} = 3$，$\bar{y} = 5$，$\hat{b} = 1.2$，则$\hat{a} = $\score{5}
\begin{enumerate}[label=\Alph*., leftmargin=2em]
\item $1.4$ \quad B. $1.6$ \quad C. $8.6$ \quad D. $11$
\end{enumerate}

\item 在残差分析中，若残差图中的点分布在一个水平的带状区域内，则说明\score{5}
\begin{enumerate}[label=\Alph*., leftmargin=2em]
\item 回归模型拟合效果不好
\item 回归模型拟合效果较好
\item 变量之间不具有线性相关关系
\item 数据存在异常值
\end{enumerate}

\item 对分类变量$X$和$Y$进行独立性检验，计算得$\chi^2 = 6.635$，则下列结论正确的是\score{5}

（参考数据：$P(\chi^2 \geqslant 3.841) = 0.05$，$P(\chi^2 \geqslant 6.635) = 0.01$）

\begin{enumerate}[label=\Alph*., leftmargin=2em]
\item 在犯错误的概率不超过$0.01$的前提下认为$X$与$Y$有关
\item 在犯错误的概率不超过$0.05$的前提下认为$X$与$Y$有关
\item 在犯错误的概率不超过$0.01$的前提下认为$X$与$Y$无关
\item 不能认为$X$与$Y$有关
\end{enumerate}

\item 已知变量$X$与$Y$的观测数据如下表：

\begin{center}
\begin{tabular}{|c|c|c|c|c|c|}
\hline
$X$ & 1 & 2 & 3 & 4 & 5 \\
\hline
$Y$ & 2 & 3 & 5 & 6 & 7 \\
\hline
\end{tabular}
\end{center}

则$Y$关于$X$的回归直线方程为\score{5}
\begin{enumerate}[label=\Alph*., leftmargin=2em]
\item $\hat{y} = 1.4x + 0.4$ \quad B. $\hat{y} = 1.3x + 0.6$ \quad C. $\hat{y} = 1.4x + 0.6$ \quad D. $\hat{y} = 1.5x + 0.2$
\end{enumerate}

\end{enumerate}

\vspace{0.5cm}

\noindent\textbf{二、多项选择题（本大题共4小题，每小题5分，共20分。每小题有多个选项符合题目要求，全部选对得5分，部分选对得2分，有选错的得0分。）}

\vspace{0.3cm}

\begin{enumerate}[label=\textbf{\arabic*.}, leftmargin=2em, resume]

\item 下列关于线性回归分析的说法正确的是\score{5}
\begin{enumerate}[label=\Alph*., leftmargin=2em]
\item 回归直线一定过样本点的中心$(\bar{x}, \bar{y})$
\item 回归直线的斜率$\hat{b}$可以为负数
\item 用最小二乘法求得的回归直线使残差平方和最小
\item 回归分析可以用来进行预测
\end{enumerate}

\item 关于相关系数$r$，下列说法正确的是\score{5}
\begin{enumerate}[label=\Alph*., leftmargin=2em]
\item $|r|$越接近1，线性相关程度越强
\item $|r|$越接近0，线性相关程度越弱
\item $r > 0$表示正相关，$r < 0$表示负相关
\item $|r| = 1$表示两个变量之间一定存在因果关系
\end{enumerate}

\item 对两个分类变量$A$和$B$的观测数据进行独立性检验，则下列说法正确的是\score{5}
\begin{enumerate}[label=\Alph*., leftmargin=2em]
\item 零假设$H_0$：$A$与$B$相互独立
\item $\chi^2$统计量的值越大，拒绝$H_0$的理由越充分
\item 当$\chi^2$的观测值大于临界值时，在相应概率下认为$A$与$B$有关
\item 独立性检验的结果是绝对正确的
\end{enumerate}

\item 已知变量$X$与$Y$的相关系数$r = -0.89$，则下列说法正确的是\score{5}
\begin{enumerate}[label=\Alph*., leftmargin=2em]
\item $X$与$Y$之间具有线性相关关系
\item $X$与$Y$之间是负相关
\item $X$增大时$Y$有减小的趋势
\item 用回归直线拟合效果一定很好
\end{enumerate}

\end{enumerate}

\vspace{0.5cm}

\noindent\textbf{三、填空题（本大题共4小题，每小题5分，共20分。）}

\vspace{0.3cm}

\begin{enumerate}[label=\textbf{\arabic*.}, leftmargin=2em, resume]

\item 已知$\bar{x} = 4$，$\bar{y} = 5$，回归直线方程为$\hat{y} = \hat{b}x + 1$，则$\hat{b} = $\hspace{3cm}。\score{5}

\item 某工厂生产的产品合格率$p$与工人工作年限$x$的相关系数$r = 0.97$，这说明两者具有\hspace{2cm}的\hspace{2cm}相关关系。\score{5}

\item 在回归分析中，决定系数$R^2$越接近\hspace{1cm}，说明回归模型的拟合效果越好。\score{5}

\item 已知根据8组数据求得$\hat{b} = 2.5$，$\hat{a} = -3$，当$x = 10$时，$y$的预测值$\hat{y} = $\hspace{3cm}。\score{5}

\end{enumerate}

\newpage

\noindent\textbf{四、解答题（本大题共6小题，共70分。解答应写出必要的文字说明、证明过程或演算步骤。）}

\vspace{0.3cm}

\begin{enumerate}[label=\textbf{\arabic*.}, leftmargin=2em, resume]

\item（10分）某兴趣小组研究学习时间（单位：小时/周）与数学成绩之间的关系，收集了6名同学的数据如下：

\begin{center}
\begin{tabular}{|c|c|c|c|c|c|c|}
\hline
学习时间$x$ & 10 & 15 & 20 & 25 & 30 & 35 \\
\hline
数学成绩$y$ & 60 & 68 & 75 & 82 & 88 & 95 \\
\hline
\end{tabular}
\end{center}

（1）求$y$关于$x$的回归直线方程$\hat{y} = \hat{b}x + \hat{a}$；

（2）预测学习时间为40小时/周的同学的数学成绩。

（参考公式：$\hat{b} = \dfrac{\displaystyle\sum_{i=1}^{n} x_i y_i - n\bar{x}\bar{y}}{\displaystyle\sum_{i=1}^{n} x_i^2 - n\bar{x}^2}$，$\hat{a} = \bar{y} - \hat{b}\bar{x}$）

\item（12分）为研究气温（$x$，单位：$^\circ\text{C}$）与冰淇淋日销量（$y$，单位：箱）的关系，某店收集了如下数据：

\begin{center}
\begin{tabular}{|c|c|c|c|c|c|}
\hline
气温$x$ & 20 & 25 & 30 & 35 & 40 \\
\hline
销量$y$ & 4 & 6 & 7 & 10 & 13 \\
\hline
\end{tabular}
\end{center}

已知$\sum_{i=1}^{5} x_i^2 = 4350$，$\sum_{i=1}^{5} x_i y_i = 1075$。

（1）求$y$关于$x$的回归直线方程；

（2）当气温为$38^\circ\text{C}$时，预测冰淇淋的日销量；

（3）求相关系数$r$，并判断线性相关程度。（参考数据：$\sum_{i=1}^{5} y_i^2 = 350$）

\item（12分）某研究机构研究了广告投入费用$x$（万元）与产品销售额$y$（万元）之间的关系，收集了8组数据，部分计算结果如下：

$\bar{x} = 5$，$\bar{y} = 8$，$\sum_{i=1}^{8} x_i^2 = 230$，$\sum_{i=1}^{8} x_i y_i = 380$，$\sum_{i=1}^{8} y_i^2 = 620$。

（1）求$y$关于$x$的回归直线方程$\hat{y} = \hat{b}x + \hat{a}$；

（2）计算决定系数$R^2$，并说明回归模型的拟合效果；

（3）当广告投入为8万元时，预测产品销售额。

\item（12分）某校对100名学生的选课意愿进行调查，结果如下：

\begin{center}
\begin{tabular}{|c|c|c|c|}
\hline
 & 选择物理 & 不选择物理 & 合计 \\
\hline
男生 & 30 & 10 & 40 \\
\hline
女生 & 20 & 40 & 60 \\
\hline
合计 & 50 & 50 & 100 \\
\hline
\end{tabular}
\end{center}

（1）根据以上数据，判断是否有$99\%$的把握认为选课意愿与性别有关；

（2）从选择物理的50名学生中随机抽取3人进行访谈，求至少有1名女生的概率。

（参考公式：$\chi^2 = \dfrac{n(ad-bc)^2}{(a+b)(c+d)(a+c)(b+d)}$，参考数据：$P(\chi^2 \geqslant 6.635) = 0.01$）

\item（12分）为研究某种植物的生长高度$y$（cm）与施肥量$x$（g）之间的关系，通过实验获得如下数据：

\begin{center}
\begin{tabular}{|c|c|c|c|c|c|c|c|}
\hline
施肥量$x$ & 1 & 2 & 3 & 4 & 5 & 6 & 7 \\
\hline
生长高度$y$ & 2.0 & 3.5 & 6.0 & 9.5 & 14 & 20 & 27 \\
\hline
\end{tabular}
\end{center}

（1）画出散点图，判断$y$与$x$之间是否具有线性相关关系；

（2）若$y$与$x$的关系更适合用指数型模型$y = ae^{bx}$描述，请通过对数变换将其转化为线性回归问题；

（3）求线性化后的回归方程，并预测当$x = 8$时$y$的值。（$\ln 2 \approx 0.693$，$\ln 3 \approx 1.099$）

\item（12分）某工厂有两条生产线生产同一产品，现从两条生产线的产品中各抽取100件进行检测，结果如下：

\begin{center}
\begin{tabular}{|c|c|c|c|}
\hline
 & 合格品 & 不合格品 & 合计 \\
\hline
生产线A & 90 & 10 & 100 \\
\hline
生产线B & 85 & 15 & 100 \\
\hline
合计 & 175 & 25 & 200 \\
\hline
\end{tabular}
\end{center}

（1）根据以上数据，判断两条生产线的产品合格率是否有显著差异（$\alpha = 0.05$）；

（2）从生产线B的100件产品中随机抽取5件进行详细检测，用$X$表示抽到的不合格品件数，求$X$的分布列和$E(X)$。

（参考数据：$P(\chi^2 \geqslant 3.841) = 0.05$，$P(\chi^2 \geqslant 6.635) = 0.01$）

\end{enumerate}

\vspace{1cm}
\begin{center}
\rule{8cm}{0.4pt}

\vspace{0.3cm}
\textbf{本试卷结束}
\end{center}

\end{document}
